\chapter{Graduate Attributes}
\label{app:outcomes}

This appendix describes how each graduate attribute (GA) was achieved through
the project, and refers to the sections of the report that contain the supporting evidence.

% =====================================================================
\section{Complex Engineering Problem}
\label{app:complex}

Designing a portable DPV potentiostat (Sections~\ref{sec:bg_and_info} and~\ref{sec:project_aim_obj_scope}) is a complex
engineering problem because first-principles were required to build solitions 
and knowledge of electrochemistry, mixed-signal
electronics and embedded software was required. The instrument must
resolve nanoampere currents while a Wi-Fi radio, display and SD card share the same supply and ground. The electrochemical cell forms part of its feedback loops, so the
behaviour of the electronics depends on what is connected to it
(Section~\ref{sec:stability}). It draws on specialist knowledge speficically regarding integrated electrochemical front-ends and the current literature
on open-source potentiostats (Chapter~\ref{chap:literature}), and on established practice for
mixed-signal design (Section~\ref{sec:hw_ground}).

The problem also shows several further characteristics of a complex engineering problem.
It involves conflicting requirements, trading sensitivity against cost, size
and battery life, within an environmental and public-health problem that has
no regulatory framework (Section~\ref{sec:dpv_applications}). There is no obvious solution to this problem, aspects such as the
architecture, the amplifier and the timing scheme each had to be chosen from
alternatives (Table~\ref{tab:afe_approaches}, Sections~\ref{sec:hstia_lptia} and~\ref{sec:sw_sampling}). The dual-interface
instrument is not offered by any surveyed design (Section~\ref{sec:selfhosted_webinterface}), and no
standard prescribes how such an instrument should be validated, so the test
method of Chapter~\ref{chap:testing} had to be devised. The work crosses into analytical
chemistry and depends on operating conditions set by an electrochemistry
research group (Section~\ref{sec:requirements}). Finally, it is a system problem, with the
front-end, power, user interface, firmware and web dashboard designed together
(Figure~\ref{fig:signal_chain},Chapters~\ref{chap:hardware_design} and~\ref{chap:software_design}).

% =====================================================================
\section{Problem Solving (GA~1)}
\label{app:ga1}

Trace-level monitoring of antiretroviral residues in South African water requires sensitive voltammetric instruments, but these are either costly and closed-source. Open-source options are often tethered to a computer or phone, so no low-cost, standalone instrument exists for DPV measurements in the field. The problem was identified in Sections~\ref{sec:bg_and_info}
and~\ref{sec:prob_statement} and formulated as an aim, scope and exclusions in Sections~\ref{sec:project_aim_obj_scope}
and~\ref{sec:scope_exclusions}. It was then turned into measurable requirements in Tables~\ref{tab:req_dpv}
and~\ref{tab:requirements_whole_sys}, so that the design could later be tested against them.

The literature was researched in Chapter~\ref{chap:literature}, first to build an understanding about electrochemical designs that depend on Section~\ref{sec:electrochemical_background} and then to survey existing instruments and locate the gap that the project fills (Section~\ref{sec:survey_designs},
Table~\ref{tab:design_comparison}). The problem was analysed from first principles rather than by
adopting a reference design. The sampling scheme follows from the decay of the
capacitive and faradaic currents (Section~\ref{sec:dpv_theory}), the stability of the
front-end from a feedback model of the cell and amplifier (Section~\ref{sec:stability}), and
the supply filtering, power chain and input filters from circuit analysis
(Sections~\ref{sec:hw_afe} to~\ref{sec:hw_power}). The conclusions made in Chapter~\ref{chap:conclusion} are based on the measurements
of Chapter~\ref{chap:testing}. When a result was unexpected, its cause was reasoned out from the data before a conclusion was drawn (Sections~\ref{sec:rtia_cal} and~\ref{sec:fc_shift}) as well as sources of deviation discussed in Section~\ref{sec:fc_deviation}.

Sustainable development was considered holistically. The problem itself is a
sustainability problem, the pollution of water by pharmaceutical residues
(Section~\ref{sec:dpv_applications}), and Chapter~\ref{sec:environmental} weighs the benefit of cheaper monitoring
against the device's own footprint. Social sciences played a small role. South Africa runs a large antiretroviral treatment programme but has no monitoring programme or regulatory limit for these residues (Section~\ref{sec:dpv_applications}).

% =====================================================================
\section{Application of Scientific and Engineering Knowledge (GA~2)}
\label{app:ga2}

GA~2 was achieved by applying mathematics, natural science, computing, engineering fundamentals and an engineering specialisation, each to a specific part of the solution. Mathematics was applied in the transfer
functions and phase-margin analysis of Section~\ref{sec:stability} and
Appendix~\ref{app:modeling}, the ferrite bead analysis (Section~\ref{sec:hw_ferrite}), the
quantisation of the DAC (Section~\ref{sec:sw_potential}), and statistics in the calibration fit and the instrument comparisons of Chapter~\ref{chap:testing}. Natural science, in the form of
electrochemistry, was applied to explain how DPV separates faradaic from
capacitive current (Sections~\ref{sec:voltammetry_family} to~\ref{sec:dpv_theory}) and to interpret the drift in the
ferricyanide results (Section~\ref{sec:fc_deviation}). Computing was applied in the firmware
of Chapter~\ref{chap:software_design} and to the digital filtering of Section~\ref{sec:sw_window}.  Engineering fundamentals and the electronic instrumentation
specialisation were applied to the amplifier, decoupling, grounding and power
design of Chapters~\ref{chap:modeling} and~\ref{chap:hardware_design}.

Mathematics, natural science and engineering science were applied in two
ways, formal analysis and modelling, and reasoning about engineering
problems. Formal analysis and modelling
appear in the cell and stability model (Table~\ref{tab:pm}) and the power-chain
estimates (Section~\ref{sec:hw_power}). Reasoning about engineering problems appears in the
interpretation of results, for example in attributing a common gain error on
all ranges to the shared conversion chain (Section~\ref{sec:rtia_cal}).

% =====================================================================
\section{Engineering Design (GA~3)}
\label{app:ga3}

GA~3 was achieved through the design of the OpenStat as a complete system to
meet the needs captured in Tables~\ref{tab:req_dpv} and~\ref{tab:requirements_whole_sys}. The creative element was expressed in
dual interface. An onboard display for standalone field use alongside a
browser dashboard, served from one sweep engine, which addresses the gap found
in Section~\ref{sec:selfhosted_webinterface} (Section~\ref{sec:sw_overview}). Designing the GUI for each interface allowed for a form of creative expression as well as the case design. The hardware is designed in Chapter~\ref{chap:hardware_design} and the firmware in Chapter~\ref{chap:software_design}. At each
major decision, the alternatives were compared before one was chosen, for
example the choice of front-end (Table~\ref{tab:afe_approaches}), the TIA (Section~\ref{sec:hstia_lptia}) and how
the scan is timed (Section~\ref{sec:hw_regulation}). The design was verified against its requirements in
Chapter~\ref{chap:testing}, and its weaknesses are acknowledged in Sections~\ref{sec:hw_regulation} and~\ref{sec:future_work}.

 Considerations made of public health and safety were that the device runs below \SI{5}{\volt}, and the lithium-polymer cell, is managed by the charger, the protection IC
and the regulator's low-voltage cut-off (Section~\ref{sec:hw_power}). Potassium ferricyanide releases cyanide in light
(Section~\ref{sec:fc_deviation}) and in contact with strong acid, so its solutions and waste
must be kept from acids and light and handled with gloves and eye protection.
Because a false negative could give false reassurance about water safety, the
OpenStat is intended as a screening tool whose positive results are confirmed
by an accredited laboratory.

Whole-life cost and net zero carbon were considered in the BOM. The component cost
(Appendix~\ref{app:bom}) is a small fraction of that of a commercial
instrument, the electrodes were reused to reduce running cost, and operating
emissions are negligible, with the power chain identified as the main
improvement (Chapter~\ref{sec:environmental}).

Resources, culture and society considerations revolved around the electrodes as they are the main consumable and their
reuse is discussed in Section~\ref{app:sustainability}. The onboard display means that no smartphone
or computer is needed, the dashboard runs in any browser, and the open-source
design allows local repair. The interface is in English only, icon-based menus
or other South African languages are a natural extension.

% =====================================================================
\section{Investigations, Experiments and Data Analysis (GA~4)}
\label{app:ga4}

GA~4 was achieved through the validation in Chapter~\ref{chap:testing}, an objective investigation where the aim was to gain insight into the performance of the OpenStat against a commercial instrument. The experiments were designed as five progressive
tests, each adding one element of a real measurement and comparing it with a
suitable reference, so that a new deviation can be traced to the element just
added (introduction to Chapter~\ref{chap:testing}, Table~\ref{tab:test_conditions}). The data were analysed
quantitatively (Tables~\ref{tab:waveform} to~\ref{tab:salt}) and interpreted with reference to the
literature where the electronics could not explain them
(Section~\ref{sec:fc_deviation}).

Critical thinking was applied in choosing how to judge the agreement of measurements. Relying on correlation alone was not accepted as evidence, the reference instrument was not treated
as the true value once its own repeatability was measured, and the difference
between instruments was judged against that repeatability
(Section~\ref{sec:fc_height}). The investigation concluded that the
instrument is not the limiting factor.

% =====================================================================
\section{Use of Engineering Tools (GA~5)}
\label{app:ga5}

GA~5 was achieved by creating, selecting, applying tools, and by stating
their limitations where they affected a result. The tools created were the
project-specific AD5941 driver and automatic range selection (Sections~\ref{sec:sw_overview}
and~\ref{sec:sw_ranging}), the firmware averaging window that replaces the on-chip notch
filter (Section~\ref{sec:sw_window}), the serial commands used for calibration
(Figure~\ref{fig:architecture}), additional Randles dummy cells (Section~\ref{sec:randles_test}) and the analysis
scripts used to compare instruments (Sections~\ref{sec:salt_test} and~\ref{sec:fc_test}). The tools selected
and applied were: KiCad and JLCPCB for the board, PlatformIO and established
libraries in C++ for the firmware (Section~\ref{sec:sw_overview}), the IPC-2221 calculator for trace width calculations 
(Section~\ref{sec:hw_trace_widths}), the laboratory instruments of Chapter~\ref{chap:testing} and \LaTeX{} for the
report. 

The limitations of each tool were recognised where they affected a result.
The accuracy of the measuring instruments limited how precisely the run time
and waveform could be verified (Sections~\ref{sec:waveform_amplitude}). Where
the datasheet was incomplete or inconsistent, conservative values were
assumed (Section~\ref{sec:stability}, Table~\ref{tab:hw_caps}). The integrated front-end limits the
compliance voltage and current ranges (Table~\ref{tab:afe_approaches}, Section~\ref{sec:sw_ranges}), and even
the reference instrument has its own scan-to-scan variation
(Section~\ref{sec:fc_height}).

% =====================================================================
\section[Professional and Technical Communication (GA~6)]{Professional and Technical\\ Communication (GA~6)}
\label{app:ga6}

GA~6 is achieved in writing through the report, which follows the structure of
the module framework and the conventions of the discipline: requirements
tables, signal-chain and flow diagrams, schematic and PCB drawings, a bill of
materials, numbered equations and IEEE referencing. Results are summarised in
tables for readers who need only the outcome (Tables~\ref{tab:testing_summary} and~\ref{tab:env_lifecycle}).

To communicate inclusively, the report explains chemical terms where they first
appear and provides lists of acronyms and definitions, since many readers will
not be electrochemists, and the summary is given in English and Afrikaans. 

% =====================================================================
\section{Individual Work (GA~8)}
\label{app:ga8}

GA~8 was achieved by carrying out every stage of the project individually,
from the literature review and requirements to the PCB design, assembly,
firmware, laboratory testing, analysis and report, with guidance from Dr T Ebrahim in weekly meetings. Time was managed with the schedule in
Appendix~\ref{app:plan_schedule}, which overlapped firmware development with board fabrication and
built in a buffer for delays.

The project crossed a disciplinary boundary into analytical chemistry,
working to operating conditions set by an electrochemistry research group
(Section~\ref{sec:requirements}) and carrying out laboratory measurements with chemical
solutions (Chapter~\ref{chap:testing}). Professional responsibility is reflected in reporting
unfavourable results and design oversights openly (Sections~\ref{sec:hw_regulation}, \ref{sec:sw_ranges}
and~\ref{sec:fc_shift}) and in stating the limits of what the instrument can claim
(Sections~\ref{sec:scope_exclusions} and~\ref{sec:testing_summary}).

% =====================================================================
\section{Independent Learning Ability (GA~9)}
\label{app:ga9}

GA~9 was achieved by learning the required knowledge the project needed independently 
outside the electronic engineering curriculum. Topics studied independently include: electrochemistry and pulse
voltammetry (Section~\ref{sec:electrochemical_background}), register-level configuration of the AD5941 from its
datasheet (Sections~\ref{sec:tia_choice} and~\ref{sec:sw_sampling}), mixed-signal grounding and decoupling
(Section~\ref{sec:hw_afe}), and measurement practice (Chapter~\ref{chap:testing}). Adaptability to new
technologies is shown by the adoption of an integrated front-end and a
self-hosted web interface (Section~\ref{app:emerging}).

From the developmental considerations listed in GA~9, there were three that were particularly impactful. Openness to feedback, feedback on content written provided a professional outlook on the work completed. Coping with
uncertainty, the ferricyanide results varied widely even on the reference
instrument, and rather than force a conclusion, the causes were sought in the
chemistry literature (Section~\ref{sec:fc_deviation}). Asking for help,  as electrochemistry lies
outside the Electrical and Electronic Engineering field, guidance was sought from Dr T Ebrahim. On reflection, what was learnt was sufficient to
design and validate the instrument on standard solutions (Sections~\ref{sec:scope_exclusions} and~\ref{sec:future_work}).
% =====================================================================
\section{Sustainable Development}
\label{app:sustainability}

The purpose of the project is itself sustainable development. A low-cost
instrument lowers the cost barrier to monitoring pharmaceutical pollution in
water (Section~\ref{sec:dpv_applications}, Chapter~\ref{sec:environmental}). Efficient resource use comes mainly from the
integrated front-end, which keeps the board small and the part count low. Environmental
impact and net zero carbon are assessed in Chapter~\ref{sec:environmental}, where operating
emissions are found to be negligible and the power chain is identified as the
main improvement. Whole-life cost is kept low by the component cost
(Appendix~\ref{app:bom}) and by updatable open-source firmware. Electrodes were cleaned and reused, but Section~\ref{sec:fc_deviation} shows the limit of
this, as ferricyanide slowly etches gold, so reuse trades material savings
against repeatability.
% =====================================================================
\section{Critical Thinking}
\label{app:critical}

Critical thinking played a crucial part in design, judging results and in setting
the scope. In design, the obvious option was not always taken, for example both
amplifiers of the AD5941 were weighed (Section~\ref{sec:hstia_lptia}). In judging results, agreement was measured against the reference instrument's own repeatability (Section~\ref{sec:fc_height}), and
systematic errors were traced from the pattern in the data rather than corrected blindly (Sections~\ref{sec:rtia_cal} and~\ref{sec:fc_shift}). In setting the scope, real water samples were excluded because a credible result would need independent
confirmation (Section~\ref{sec:scope_exclusions}).

% =====================================================================
\section{Emerging Technologies}
\label{app:emerging}

Three emerging technologies were evaluated. Integrated electrochemical
front-ends were compared with discrete designs in Table~\ref{tab:afe_approaches} and adopted, and their limitations were then evaluated (Sections~\ref{sec:sw_ranges} and~\ref{sec:rtia_cal}). Web-hosted instruments were compared with smartphone-app designs in Sections~\ref{sec:early_opensource} and~\ref{sec:selfhosted_webinterface}; the web-hosted approach was adopted and extended with an onboard display.